\section{Fazit und Ausblick}
\subsection{Fazit}
Im Rahmen dieser Studienarbeit wurde das automatisierte Modell-Hochregallager der
Dualen Hochschule Baden-Württemberg gezielt weiterentwickelt und um wesentliche
funktionale sowie sicherheitstechnische Aspekte ergänzt. Ziel war es, bestehende Defizi
te im Bereich der Sicherheit, der Programmstruktur und der Prozesslogik zu analysieren
und durch geeignete technische Maßnahmen zu beheben. Die durchgeführten Arbei
ten leisten einen wichtigen Beitrag zur funktionalen Erweiterung und zur praxisnahen
Ausgestaltung des Systems.
Ein zentraler Schwerpunkt der Arbeit lag auf der Integration einer berührungslos wir
kenden Schutzeinrichtung in Form eines Lichtgatters. Durch die Implementierung der
Lichtschranke in Verbindung mit einem IO-Link-fähigen Sensor konnte eine zuverläs
sige Überwachung des Gefahrenbereichs realisiert werden. Die detaillierte Analyse
der verfügbaren Prozesswerte sowie die gezielte Auswahl des Modus „Number of Be
ams Occupied (NBO)“ ermöglichen eine klare Unterscheidung zwischen zulässigen
Prozessbewegungen, wie dem Einfahren des Zylinders, und unzulässigen Eingriffen
in den Arbeitsbereich. Dadurch wird eine situationsabhängige und zugleich robuste
Sicherheitsabschaltung gewährleistet.
Die sicherheitstechnische Logik wurde vollständig in die bestehende SPS-Architektur
integriert und über einen eigenen Funktionsbaustein realisiert. Durch die zentrale
Freigabevariable wird sichergestellt, dass sämtliche Aktoren des Systems nur dann
angesteuert werden, wenn keine Störung vorliegt und der Gefahrenbereich frei ist.
Dieses Vorgehen erhöht nicht nur die Sicherheit des Systems, sondern verbessert
auch die Nachvollziehbarkeit und Wartbarkeit der Steuerungssoftware. Die Kopplung
der Sicherheitsfunktion mit dem vorhandenen Störungsmanagement stellt sicher, dass
sicherheitsrelevante Zustände eindeutig erkannt, verarbeitet und visualisiert werden.
Ein weiterer wesentlicher Bestandteil der Arbeit war die Einführung eines Umlagerungs
prozesses. Im Gegensatz zu den bereits vorhandenen Ein- und Auslagerungsfunktionen
handelt es sich hierbei um einen rein internen logistischen Prozess, der zur Optimie
rung der Lagerstruktur dient. Die Entwicklung einer eigenständigen Schrittkette erwies
sich als notwendig, um die zusätzlichen Anforderungen hinsichtlich Bewegungslogik,
Verzweigungen und Prozesssicherheit abzubilden. Durch die Kombination bestehender
Grundlogiken mit neuen Entscheidungsstrukturen konnte ein flexibler und nachvollzieh
barer Umlagerungsablauf realisiert werden.
Darüber hinaus zeigt die Arbeit deutlich die Vorteile der textbasierten Programmierung
in Structured Text (ST) gegenüber grafischen Programmiersprachen wie FUP. Insbeson
dere bei komplexeren Abläufen, sicherheitsrelevanten Abfragen und mathematischen
Vergleichen bietet ST eine deutlich höhere Übersichtlichkeit und Skalierbarkeit. Die
Umstellung einzelner Funktionsbereiche auf ST trägt langfristig zu einer besseren Les
barkeit, einfacheren Fehlersuche und verbesserten Erweiterbarkeit des Programms
bei.
Zusammenfassend lässt sich festhalten, dass die in dieser Studienarbeit umgesetzten
Maßnahmen die funktionale Sicherheit, Prozessstabilität und didaktische Qualität des
Modell-Hochregallagers deutlich verbessern. Das System bildet nun nicht nur grund
legende logistische Prozesse ab, sondern berücksichtigt auch sicherheitstechnische
Anforderungen, wie sie in realen industriellen Anlagen zwingend erforderlich sind. Da
mit stellt das Hochregallager eine praxisnahe und zukunftsfähige Lernplattform für
Studierende im Bereich der Automatisierungstechnik dar

\subsection{Ausblick}

Die im Rahmen dieser Studienarbeit umgesetzten Erweiterungen bilden eine stabile
technische und konzeptionelle Grundlage für die weitere Entwicklung des automatisier
ten Modell-Hochregallagers. Insbesondere durch die Integration der sicherheitsgerichte
ten Lichtschranke, die Einführung des Umlagerungsprozesses sowie die zunehmende
Nutzung der Programmiersprache Structured Text wurde das System sowohl funktional
als auch strukturell deutlich verbessert. Aufbauend auf diesen Ergebnissen wird sich die
darauffolgende Arbeit auf die Implementierung eines auftragsbasierten Materialflusses
konzentrieren.
Der zentrale Fokus der nächsten Entwicklungsstufe liegt auf der Einführung eines
Auftragsmanagements mithilfe eines QR-Code-Readers. Ziel ist es, Ein- und Auslager
prozesse nicht mehr ausschließlich manuell oder fest vorgegeben auszulösen, sondern
diese auf Basis eindeutig definierter Aufträge automatisch zu steuern. Hierdurch soll
das Hochregallager stärker an reale industrielle Anwendungen angelehnt werden, in
denen Materialbewegungen in der Regel durch digitale Aufträge aus übergeordneten
Systemen angestoßen werden.
Im Rahmen der nächsten Arbeit soll ein QR-Reader als zusätzliche Identifikations- und
Eingabeschnittstelle in das bestehende System integriert werden. Über den QR-Code
sollen auftragsrelevante Informationen wie Lagerplatz, Material-ID, Prozessart (Einla
gern, Auslagern oder Umlagern) sowie gegebenenfalls Prioritäten oder Zielpositionen
erfasst werden. Die ausgelesenen Daten werden anschließend von der SPS verar
beitet und in strukturierter Form an die Prozesssteuerung übergeben. Damit entsteht
eine direkte Verbindung zwischen physischer Auftragserfassung und automatisierter
Anlagensteuerung.
Ein wesentlicher Bestandteil des Auftragsmanagements wird die Definition geeigneter
Datenstrukturen innerhalb der SPS sein. Hierzu zählen unter anderem Auftragsobjekte,
Warteschlangen sowie Statusinformationen zur Verfolgung des aktuellen Bearbeitungs
stands. Durch die Verwendung strukturierter Datentypen in Structured Text können
Aufträge klar definiert, verwaltet und sequenziell abgearbeitet werden. Dies ermöglicht
eine saubere Trennung zwischen Auftragsebene und Bewegungs- bzw. Sicherheitslogik
und erhöht die Modularität des Gesamtsystems.
Darüber hinaus bietet ein auftragsbasiertes System die Möglichkeit, bestehende Pro
zesse wie Einlagerung, Auslagerung und Umlagerung in einer übergeordneten Logik
zusammenzuführen. Anstatt einzelne Funktionen separat auszulösen, kann der Anla
genbetrieb künftig vollständig über Aufträge gesteuert werden. Dies eröffnet Potenzial
für eine automatische Priorisierung von Aufträgen, eine flexible Reihenfolge der Abarbei
tung sowie eine bessere Auslastung der Anlage. Insbesondere in Kombination mit dem
bereits implementierten Umlagerungsprozess lassen sich so optimierte Lagerstrategien
realisieren.
Ein weiterer Aspekt der kommenden Arbeit betrifft die Erweiterung der Bedien- und
Visualisierungsebene. Die über den QR-Reader eingelesenen Aufträge sollen auf
dem HMI übersichtlich dargestellt werden, inklusive Informationen zum aktuellen Auf
tragsstatus, zur Restbearbeitungszeit und zu möglichen Störungen. Dadurch wird die
Transparenz für den Bediener erhöht und der Anlagenzustand jederzeit nachvollziehbar
dargestellt. Ergänzend könnte eine Historie abgeschlossener Aufträge implementiert
werden, um Prozessabläufe analysieren und optimieren zu können.
Auch im Hinblick auf die Sicherheit ergeben sich neue Anforderungen. Das Auftrags
management muss eng mit den bestehenden Sicherheitsfunktionen verknüpft werden,
sodass Aufträge bei aktiven Störungen oder unterbrochenen Schutzfunktionen nicht
gestartet oder automatisch pausiert werden. Dadurch wird sichergestellt, dass der
auftragsbasierte Betrieb jederzeit den sicherheitstechnischen Anforderungen entspricht
und keine gefährlichen Zustände entstehen.
Zusammenfassend stellt die Implementierung eines QR-basierten Auftragsmanage
ments einen konsequenten nächsten Schritt in der Weiterentwicklung des automatisier
ten Modell-Hochregallagers dar. Sie ermöglicht eine realitätsnahe Abbildung moderner
intralogistischer Prozesse und schafft die Grundlage für eine flexible, skalierbare und
industrieorientierte Anlagensteuerung. Die Ergebnisse der vorliegenden Arbeit bilden
hierfür eine solide Basis und ermöglichen es, das System in der folgenden Theoriephase
gezielt um eine zentrale Funktion moderner Lagerlogistik zu erweitern

